\chapterwithopening{Literature Review and Theoretical Framework}{chap:literature_review}{
  \chaptercontentsitem{sec:literature_review_introduction}
  \chaptercontentsitem{sec:evolution_of_digital_education_platforms}
  \chaptercontentsitem{sec:research_gap_integrated_ecosystem}
  \chaptercontentsitem{sec:artificial_intelligence_in_education}
  \chaptercontentsitem{sec:benchmarking_fragmented_vs_integrated}
  \chaptercontentsitem{sec:theoretical_technical_foundations_eduverse}
}

\section{Introduction}\label{sec:literature_review_introduction}

This literature review establishes the conceptual and technical foundations
relevant to EduVerse and situates the project within existing approaches to
digital education. Because the platform combines several forms of academic
activity, its theoretical basis extends across digital learning systems,
communication and synchronous teaching, artificial intelligence in education,
access control, mobile learning, and extensible platform design.

Together, these areas provide the basis for examining how educational
activities can remain coordinated while supporting different users,
institutional requirements, and modes of participation. The review also
identifies limitations in fragmented educational environments and the design
principles that motivate a more connected approach.

\section{Evolution of Digital Education Platforms}
\label{sec:evolution_of_digital_education_platforms}

\subsection{From Traditional Classroom Management to Digital Learning Systems}
\label{subsec:traditional_to_digital_learning_systems}

Educational management has gradually moved from paper-based and face-to-face
coordination toward persistent digital systems capable of organizing materials,
communication, submissions, and academic records at institutional scale.
Learning management systems emerged as part of this transition, providing
shared digital spaces through which teaching and learning could be organized
more consistently \cite{chugh2023edtech,bygstad2022digital}.

This transition also changed user expectations. Students and teachers
increasingly expect continuous access, asynchronous participation, traceable
communication, and structured academic follow-up rather than isolated
interactions \cite{chugh2023edtech}. Modern educational platforms are therefore
evaluated not only by their ability to store resources, but also by how
effectively they support ongoing academic activity.

\subsection{Limitations of Fragmented Learning Management Tools}
\label{subsec:limitations_of_fragmented_learning_management_tools}

Research on educational digitalization shows that institutions commonly rely
on combinations of specialized systems for learning management, communication,
synchronous teaching, collaboration, and practical work. Although individual
tools may perform their intended functions effectively, using multiple
disconnected environments can weaken continuity across academic activities and
increase coordination overhead \cite{bygstad2022digital,oecd2023outlook}.

Such fragmentation can require users to repeatedly move between separate
contexts while institutions manage permissions, visibility, communication, and
responsibility across different services. The resulting challenge is
maintaining consistent academic coordination when related activities are
distributed across separate tools
\cite{bygstad2022digital,oecd2023outlook}.

\section{Research Gap and Integrated Learning Ecosystems}
\label{sec:research_gap_integrated_ecosystem}

The gap relevant to EduVerse concerns coordinated operation across individual
educational tools. Existing systems can support communication,
learning resources, assessment, synchronous teaching, collaboration, and
practical work independently; however, institutions also require these
activities to operate within consistent boundaries of responsibility, access,
and academic context.

This creates a broader requirement for educational environments that combine
functional integration with institutional flexibility. Integration alone may be insufficient when different institutions or classes are expected to use the same capabilities in the same way. Conversely, flexibility becomes difficult to
manage when related workflows are distributed across unrelated systems.

EduVerse addresses this gap by adopting an integrated but configurable
educational model. The significance of this approach lies not in introducing
entirely new categories of educational tools, but in coordinating familiar
academic functions within a shared institutional context.

\section{Artificial Intelligence in Education}
\label{sec:artificial_intelligence_in_education}

Artificial intelligence has become increasingly relevant to education because
it can support explanation, summarization, drafting, and personalized academic
assistance when used within appropriate pedagogical boundaries
\cite{bond2024ai,labadze2023chatbots}. Its educational value, however, depends
not only on its ability to generate content, but also on how its use affects
teacher judgment, academic integrity, and the learner's engagement with source
material.

For EduVerse, this literature supports treating AI as a supplementary learning
capability rather than as a replacement for instructional authority or
established academic resources.

\subsection{Context-Aware Learning Assistance}
\label{subsec:context_aware_learning_assistance}

AI assistance is generally more useful when it is informed by the learner's
current academic context. Responses related to the active class, learning
material, or task can provide more relevant support than generic answers
detached from the educational setting
\cite{alsafari2024teaching,labadze2023chatbots}. Context awareness can therefore
improve both relevance and pedagogical alignment.

This principle supports the use of embedded AI assistance in EduVerse, where
learning support is intended to remain connected to the academic environment in
which it is requested.

\subsection{Material Summarization and Learning Support}
\label{subsec:material_summarization_and_learning_support}

Long or numerous educational resources can make it difficult for students to
quickly recover the structure and key ideas of a topic. AI-supported
summarization can assist by foregrounding important concepts and reducing the
effort required to re-enter complex material \cite{xie2025summarization}.

Such summaries should remain supportive rather than substitutive. In the
context of EduVerse, summarization is therefore treated as an aid for navigating
learning material rather than as a replacement for the original resource.

\subsection{Assessment and Feedback Support}
\label{subsec:assessment_and_feedback_support}

The literature indicates that AI can contribute to formative learning support,
drafting, and feedback-related activities when used within clear educational
boundaries \cite{labadze2023chatbots}. At the same time, unrestricted use can
weaken authentic assessment and blur the distinction between assistance and
task completion \cite{unesco2023genai,labadze2023chatbots}.

These concerns support a cautious role for AI within EduVerse. AI-assisted
learning should complement teacher-led assessment and study activities without
replacing instructional authority or undermining the integrity of assessed
work.

\section{Benchmarking: Fragmented Tools vs. Integrated EduVerse Approach}
\label{sec:benchmarking_fragmented_vs_integrated}

The purpose of this benchmarking comparison is not to rank individual products,
but to contrast two different organizational models for digital learning:
workflows distributed across specialized tools and workflows coordinated within
a shared educational environment.
Table~\ref{tab:fragmented_vs_integrated_eduverse} summarizes the main structural
differences relevant to EduVerse.

\begingroup
\footnotesize
\setlength{\tabcolsep}{4pt}
\renewcommand{\arraystretch}{1.15}
\begin{longtable}{|>{\raggedright\arraybackslash}p{0.17\textwidth}|>{\raggedright\arraybackslash}p{0.34\textwidth}|>{\raggedright\arraybackslash}p{0.34\textwidth}|}
\caption{Fragmented Educational Tools vs. Integrated EduVerse Approach}
\label{tab:fragmented_vs_integrated_eduverse} \\
\hline
\rowcolor{gray!15}
\textbf{Workflow Area} &
\textbf{Fragmented Reference Model} &
\textbf{Integrated EduVerse Approach} \\
\hline
\endfirsthead
\hline
\rowcolor{gray!15}
\textbf{Workflow Area} &
\textbf{Fragmented Reference Model} &
\textbf{Integrated EduVerse Approach} \\
\hline
\endhead
\hline
\multicolumn{3}{r}{\small\textit{Continued on the next page}} \\
\endfoot
\hline
\endlastfoot
Learning Content and Communication &
Learning resources, announcements, communication, and related class activity
may be distributed across different systems, requiring users to move repeatedly
between contexts. &
Learning content and class communication remain connected to the same
organizational and class context. \\
\hline
Synchronous Collaboration &
Live teaching, presentation interaction, and collaborative whiteboard activity
may depend on separate communication and collaboration tools. &
Synchronous learning and collaborative activity remain associated with the
surrounding class environment. \\
\hline
Assessment and Practical Learning &
Assignments, examinations, results, and practical activities such as programming
may be handled through separate academic or specialist systems. &
Assessment and supported practical learning activities can remain connected to
the wider academic workflow. \\
\hline
Governance and Configurability &
Roles, permissions, visibility, and available capabilities may need to be
managed separately across different services. &
Institutional governance and supported feature configuration are coordinated
within the same educational environment. \\
\hline
Integration and Extension &
Extending a distributed environment may require additional integration between
independent systems and duplicated coordination logic. &
A shared platform model provides a common foundation for supported extensions
and connected clients. \\
\hline
\end{longtable}
\endgroup

The comparison indicates that the main distinction lies in coordination across
related academic activities rather than in any single feature. EduVerse
therefore draws its value from the relationship between these capabilities and
the institutional context in which they operate.

\section{Theoretical and Technical Foundations of EduVerse}
\label{sec:theoretical_technical_foundations_eduverse}

The preceding literature identifies several principles relevant to the design
of an integrated educational environment. The following subsections summarize
the foundations most directly related to EduVerse, including web-based
educational systems, mobile companion applications, synchronous learning,
integrated programming environments, extensibility, access control, and data
management.

\subsection{Web-Based Educational Platforms}
\label{subsec:web_based_educational_platforms}

Web platforms remain central to comprehensive educational systems because they
can support rich administrative interfaces, detailed authoring workflows, and
persistent academic records. Their value lies not only in browser accessibility,
but also in their capacity to accommodate complex institutional and
instructional activities within a consistent interface.

These characteristics support the use of the web as the primary client for
comprehensive EduVerse workflows, while more focused forms of participation can
be extended through other clients.

\subsection{Mobile Companion Learning Applications}
\label{subsec:mobile_companion_learning_applications}

Mobile educational applications are often most effective when they support
frequent and lightweight participation rather than reproduce the full
complexity of a comprehensive academic platform. Notifications, quick updates,
access to materials, communication, and entry into live activities are examples
of interactions that align naturally with mobile usage
\cite{rangel2023mobile}.

This literature supports the companion-client model adopted by EduVerse, in
which mobile access extends selected educational interactions while the web
application retains broader functional depth.

\subsection{Real-Time Virtual Classrooms}
\label{subsec:real_time_virtual_classrooms}

Synchronous online teaching remains important because it supports immediacy,
responsiveness, and a stronger sense of instructional presence
\cite{lee2024synchronous}. Its educational effectiveness, however, is
influenced by how well the live activity remains connected to the wider
learning context.

For EduVerse, this principle supports embedding synchronous participation
within the class environment rather than treating a live session as an isolated
event disconnected from communication, materials, and subsequent academic
activity.

\subsection{Integrated Coding Environments for Programming Education}
\label{subsec:integrated_coding_environments_programming_education}

Programming education illustrates the cost of fragmented digital workflows
particularly clearly. When explanation and practical coding are separated
across unrelated environments, students must repeatedly switch between learning
and execution contexts. Integrated coding environments can reduce this
separation by bringing practical work closer to the surrounding instructional
activity \cite{noor2023ide}.

This principle supports the inclusion of browser-based coding within EduVerse
as a way of maintaining a closer relationship between programming practice and
the wider learning context.

\subsection{Extensible Educational Ecosystems and Governed Openness}
\label{subsec:extensible_ecosystems_governed_openness}

Educational institutions rarely require every available capability in every
learning context. Modular platform approaches address this variation by
allowing functionality to be combined without forcing a rigid all-or-nothing
structure \cite{sorkun2022modularity}. Extensibility is most useful, however,
when it remains subject to clear institutional control.

The same principle applies to openness. Broader access to learning resources or
participation can be valuable when visibility and authority remain governed by
institutional rules \cite{oecd2023outlook}. These principles support EduVerse's
configurable approach without requiring every organization or class to operate
identically.

\subsection{Security, Roles, and Data Management}
\label{subsec:security_roles_and_data_management}

Role management is a foundational requirement in educational systems because
students, teachers, and administrators interact with academic information from
different positions of responsibility. Role-Based Access Control provides an
established approach for separating permissions and responsibilities among
different actors \cite{kabier2023rbac}.

This principle supports the role-aware access model used by EduVerse, where
academic and administrative actions are differentiated according to the user's
institutional context.

Data management is equally important. Educational systems must preserve
academic records, protect sensitive information, and provide structured access
to learning materials and results \cite{liu2023privacy,kabier2023rbac}. These
requirements provide the theoretical basis for treating data protection,
controlled access, and academic-history preservation as core system concerns.

\enlargethispage{3\baselineskip}
Together, these foundations establish the academic and technical rationale for
EduVerse as an integrated educational platform. The literature indicates that
effective digital learning environments depend not only on individual
capabilities, but also on how communication, participation, assessment, access
control, and learning support are coordinated. The next chapter builds on these
foundations by presenting the methodology used to analyze, design, implement,
and refine the system.
